% !TeX encoding = UTF-8
\documentclass[a4paper,11pt]{article}

\usepackage{fontspec}
\usepackage{polyglossia}
\setdefaultlanguage{russian}
\setmainfont{Times New Roman}

\usepackage{amsmath,amssymb,mathtools}
\usepackage{geometry}
\usepackage{enumitem}
\usepackage{hyperref}
\usepackage{microtype}

\geometry{left=1.7cm,right=1.7cm,top=1.6cm,bottom=1.7cm}
\hypersetup{colorlinks=true,linkcolor=black,urlcolor=black}
\setlist[itemize]{noitemsep,topsep=2pt,leftmargin=1.2em}
\setlist[enumerate]{noitemsep,topsep=2pt,leftmargin=1.4em}
\setlength{\parindent}{0pt}
\setlength{\parskip}{4pt}
\sloppy

\newcommand{\R}{\mathbb{R}}
\newcommand{\dd}{\,d}
\newcommand{\argmax}{\operatorname*{arg\,max}}
\newcommand{\argmin}{\operatorname*{arg\,min}}
\newcommand{\ticket}[2]{%
  \clearpage
  \section*{Билет #1. #2}
  \addcontentsline{toc}{section}{Билет #1. #2}
}

\begin{document}

\begin{titlepage}
\centering
{\Large Московский государственный университет имени М.В. Ломоносова\par}
\vspace{0.4cm}
{\large Факультет вычислительной математики и кибернетики\par}
\vspace{0.4cm}
{\large Кафедра оптимального управления\par}
\vfill
{\LARGE\bfseries Экзаменационные ответы к государственному экзамену\par}
\vspace{0.5cm}
{\large Магистерская программа\\
«Математические методы моделирования и методы оптимизации управляемых процессов»\par}
\vfill
{\large Краткая версия для ответа на экзамене\par}
\vspace{0.5cm}
{\large Москва, 2026\par}
\end{titlepage}

\tableofcontents

\ticket{1}{Обратные задачи для систем ОДУ. Метод динамической регуляризации}

Рассматривается динамическая система
\[
 \dot y(t)=f(t,y(t))u(t)+g(t,y(t)),\quad y(0)=y_0,\quad t\in[0,T],
\]
где состояние \(y(t)\in\R^n\), управление \(u(t)\in P\subset\R^r\), а функции \(f,g\), начальное состояние и горизонт \(T\) известны.
Прямая задача состоит в нахождении траектории \(y(\cdot)\) по заданному управлению \(u(\cdot)\).
Обратная задача состоит в восстановлении управления \(u(\cdot)\) по наблюдаемой траектории \(y(\cdot)\).

Обратная задача, вообще говоря, некорректна: решение может быть неединственным и неустойчивым относительно ошибок наблюдения.
Поэтому выбирают нормальное управление, то есть управление минимальной \(L_2\)-нормы среди всех управлений, порождающих наблюдаемую траекторию:
\[
 U_*=\{u\in U:\ y(t,u)=y^r(t)\ \text{на }[0,T]\},\qquad
 u_*=\argmin_{u\in U_*}\|u\|_{L_2(0,T)}.
\]
При липшицевости \(f,g\) и компактности \(P\) прямая задача имеет единственное решение, а при выпуклости \(P\) множество \(U_*\) выпукло и замкнуто; поэтому нормальное управление существует и единственно.

На практике известна не точная траектория, а наблюдения
\[
 \|y(t_i)-\tilde y_i\|\le \delta,\quad 0=t_0<t_1<\dots<t_N=T,\quad h=\max_i(t_{i+1}-t_i).
\]
Метод динамической регуляризации Осипова--Кряжимского строит кусочно-постоянное приближение \(\tilde u_{\delta,h}\).
Вводится траектория-поводырь \(z_i\):
\[
 \frac{z_{i+1}-z_i}{h}=f(t_i,\tilde y_i)u_i+g(t_i,\tilde y_i),\qquad z_0=\tilde y_0.
\]
На каждом шаге управление выбирается из задачи регуляризованной минимизации
\[
 T_i(v)=2\langle z_i-\tilde y_i,\ f(t_i,\tilde y_i)v\rangle+\alpha \|v\|^2
 \to \min_{v\in P},
\]
точнее \(u_i\in P\) берется так, что \(T_i(u_i)\le \inf_{v\in P}T_i(v)+\varepsilon\).
Стабилизатор Тихонова \(\alpha\|v\|^2\) обеспечивает устойчивость.

Основной результат: при согласованном стремлении параметров к нулю
\[
 \alpha,\delta,\varepsilon,h\to0,\qquad \frac{\delta+h+\varepsilon}{\alpha}\to0,
\]
траектория поводыря сходится к точной траектории, а построенные управления сходятся к нормальному управлению:
\[
 \|z_h-y\|_{C[0,T]}\to0,\qquad
 \|\tilde u_{\delta,h}-u_*\|_{L_2(0,T)}\to0.
\]
Итак, метод динамической регуляризации заменяет некорректное восстановление управления устойчивой пошаговой процедурой с регуляризацией.

\ticket{2}{Обратные задачи для параболических уравнений. Метод динамической регуляризации}

Типичная модель -- уравнение теплопроводности с распределенным управлением:
\[
\begin{cases}
 y_t(t,x)=a^2y_{xx}(t,x)+u(t,x), & (t,x)\in(0,T)\times(0,l),\\
 y(t,0)=y(t,l)=0,\\
 y(0,x)=y_0(x).
\end{cases}
\]
Прямая задача: по \(u(t,x)\) найти \(y(t,x)\).
Обратная задача: по наблюдениям \(y(t,x)\) восстановить правую часть \(u(t,x)\).
Управления рассматриваются в пространстве \(L_2(Q)\), \(Q=(0,T)\times(0,l)\).

Переходят к абстрактной эволюционной форме. Пусть
\[
 H=L_2(0,l),\qquad V=H_0^1(0,l),\qquad V\subset H\simeq H^*\subset V^*,
\]
а оператор \(A:V\to V^*\) задается формулой
\[
 \langle Ay,v\rangle=a^2\int_0^l y_x(x)v_x(x)\dd x.
\]
Тогда уравнение записывается как
\[
 \dot y(t)+Ay(t)=u(t),\qquad y(0)=y_0.
\]
Оператор \(A\) симметричен, положительно определен, а вложение \(V\subset H\) компактно. Эти свойства обеспечивают корректность прямой задачи.

Наблюдения заданы в дискретные моменты времени:
\[
 \|y(t_i)-\tilde y_i\|_{L_2(0,l)}\le \delta.
\]
Как и в случае ОДУ, строится траектория-поводырь \(z_h\), удовлетворяющая дискретной операторной схеме
\[
 \dot z_h(t)+Az_h(t)=u_i,\qquad t\in(t_i,t_{i+1}).
\]
Управление на очередном промежутке выбирается из регуляризованной задачи:
\[
 \varphi_i(v)=2\langle z_h(t_i)-\tilde y_i,\ v\rangle_{L_2}
 +\alpha\|v\|_{L_2}^2\to \min_{v\in P},
\]
где \(P\subset L_2(0,l)\) -- выпуклое, замкнутое и ограниченное множество допустимых значений управления.

Если \(u_i\) выбирается с точностью \(\varepsilon\), то получаем кусочно-постоянное управление
\[
 \tilde u(t)=u_i,\qquad t\in(t_i,t_{i+1}).
\]
Теорема сходимости утверждает, что при \(\delta,\varepsilon,\alpha,h\to0\)
\[
 \|z_h-y\|_{C([0,T];L_2(0,l))}\to0,
\]
а при дополнительном согласовании
\[
 \frac{\sqrt h+\delta+\varepsilon}{\alpha}\to0
\]
имеет место сходимость управления к нормальному решению:
\[
 \|\tilde u-u_*\|_{L_2((0,T);L_2(0,l))}\to0.
\]
Смысл метода: неизвестная правая часть теплового уравнения восстанавливается не прямым дифференцированием зашумленных данных, а устойчивым управлением поводыря, близкого к наблюдаемому процессу.

\ticket{3}{Обратные задачи для гиперболических уравнений. Метод динамической регуляризации}

Для гиперболического случая основной пример -- волновое уравнение:
\[
\begin{cases}
 y_{tt}(t,x)=a^2y_{xx}(t,x)+u(t,x), &(t,x)\in(0,T)\times(0,l),\\
 y(t,0)=y(t,l)=0,\\
 y(0,x)=y_0(x),\quad y_t(0,x)=y_1(x).
\end{cases}
\]
Обратная задача состоит в восстановлении \(u(t,x)\) по наблюдениям состояния или скорости. В отличие от параболического случая система второго порядка по времени, поэтому естественно рассматривать пару \((y,\dot y)\).

Вводятся пространства
\[
 V=H_0^1(0,l),\qquad H=L_2(0,l),\qquad V\subset H\simeq H^*\subset V^*,
\]
и положительный симметрический оператор \(A:V\to V^*\),
\[
 \langle Ay,v\rangle=a^2\int_0^l y_xv_x\dd x.
\]
Тогда волновое уравнение имеет абстрактный вид
\[
 \ddot y(t)+Ay(t)=u(t),\qquad y(0)=y_0,\quad \dot y(0)=y_1.
\]
Допустимые управления принадлежат \(L_2((0,T);H)\) и почти всюду принимают значения в выпуклом замкнутом ограниченном множестве \(P\subset H\).

Метод динамической регуляризации строится аналогично параболическому случаю, но поводырь учитывает вторую производную по времени:
\[
 \ddot z_h(t)+Az_h(t)=u_i,\qquad t\in(t_i,t_{i+1}).
\]
На каждом шаге используется текущая информация о наблюдаемой траектории и выбирается управление из задачи
\[
 2\langle \dot z_h(t_i)-\tilde y_i,\ v\rangle_H+\alpha\|v\|_H^2
 \to \min_{v\in P},
\]
или из эквивалентной формы, зависящей от того, наблюдается состояние или скорость.
Регуляризующий параметр \(\alpha>0\) подавляет неустойчивость, возникающую при попытке восстановить \(u\) через вторую производную зашумленных данных.

Нормальное управление определяется так же:
\[
 u_*=\argmin\{\|u\|_{L_2((0,T);H)}:\ y(\cdot,u)=y^r(\cdot)\}.
\]
При стандартных условиях на оператор \(A\), множество \(P\) и согласовании параметров дискретизации, ошибки и регуляризации выполняется
\[
 z_h\to y,\qquad \tilde u_{\delta,h}\to u_*
\]
в соответствующих энергетических нормах.

Главная идея ответа: для гиперболического уравнения обратная задача также некорректна, но метод динамической регуляризации строит устойчивое приближение неизвестной характеристики через вспомогательную динамическую систему-поводырь и регуляризованную пошаговую минимизацию.

\ticket{4}{Одномерная модель Рамсея на бесконечном горизонте. Особый режим}

Рассматривается одномерная модель оптимального экономического роста:
\[
 \dot x(t)=u(t)Ax^\varepsilon(t)-\mu x(t),\quad x(0)=x_0>0,\quad u(t)\in[0,1],
\]
\[
 J(u)=\int_0^\infty e^{-\nu t}(1-u(t))x^\varepsilon(t)\dd t\to\max.
\]
Здесь \(x\) -- удельный капитал, \(u\) -- доля выпуска, направляемая в инвестиции, \(1-u\) -- доля потребления, \(\mu\) -- амортизация, \(\nu\) -- коэффициент дисконтирования.

Гамильтониан:
\[
 H=e^{-\nu t}(1-u)x^\varepsilon+\psi(ux^\varepsilon-\mu x).
\]
Условие максимума по \(u\in[0,1]\) дает переключающую функцию
\[
 \Phi(t)=\psi(t)x^\varepsilon(t)-e^{-\nu t}x^\varepsilon(t).
\]
Если \(\Phi>0\), то \(u=1\); если \(\Phi<0\), то \(u=0\); если \(\Phi\equiv0\) на интервале, возникает особый режим.

В часто используемом случае \(A=1,\ \varepsilon=\frac12\):
\[
 \dot x=u\sqrt{x}-\mu x,\qquad
 J=\int_0^\infty e^{-\nu t}(1-u)\sqrt{x}\dd t\to\max.
\]
Особое состояние определяется условиями стационарности и имеет вид
\[
 x_{\mathrm{s}}=\frac{1}{4(\mu+\nu)^2},\qquad
 u_{\mathrm{s}}=\mu\sqrt{x_{\mathrm{s}}}=\frac{\mu}{2(\mu+\nu)}\in(0,1).
\]
На особом режиме капитал постоянен:
\[
 \dot x=0,\qquad x(t)\equiv x_{\mathrm{s}},\qquad u(t)\equiv u_{\mathrm{s}}.
\]

Качественная структура оптимального решения:
\begin{itemize}
 \item если \(x_0<x_{\mathrm{s}}\), то сначала \(u=1\), капитал растет до \(x_{\mathrm{s}}\), затем начинается особый режим;
 \item если \(x_0=x_{\mathrm{s}}\), то сразу реализуется особый режим;
 \item если \(x_0>x_{\mathrm{s}}\), то сначала \(u=0\), капитал снижается до \(x_{\mathrm{s}}\), затем начинается особый режим.
\end{itemize}
Таким образом, оптимальное управление имеет не более одной точки переключения: сначала применяется крайнее управление \(0\) или \(1\), затем система выходит на бесконечный особый участок.
Экономический смысл: оптимальная политика выводит экономику на «золотой» уровень капитала, при котором поддерживается наилучший баланс между инвестициями и текущим потреблением.

\ticket{5}{Двумерная задача распределения ресурсов в двухсекторной модели Кобба--Дугласа}

Рассматривается модель двухсекторной экономики:
\[
\begin{cases}
 \dot x_1=uF(x),\quad x_1(0)=x_{10}>0,\\
 \dot x_2=(1-u)F(x),\quad x_2(0)=x_{20}>0,\\
 J=x_2(T)\to\max,\qquad u(t)\in[0,1],
\end{cases}
\]
где \(F(x)=Ax_1^{\varepsilon_1}x_2^{\varepsilon_2}\), \(\varepsilon_1,\varepsilon_2>0\), \(\varepsilon_1+\varepsilon_2=1\).
Управление \(u\) задает долю выпуска, направленную в первый сектор; цель -- максимизировать конечный уровень второго сектора.

После замены переменных задача приводится к виду
\[
\begin{cases}
 \dot y_1=\dfrac{u}{\varepsilon_1}F(y),\\[4pt]
 \dot y_2=\dfrac{1-u}{\varepsilon_2}F(y),\\
 J=y_2(T)\to\max.
\end{cases}
\]
Гамильтониан:
\[
 H=F(y)\left[\left(\frac{\psi_1}{\varepsilon_1}-\frac{\psi_2}{\varepsilon_2}\right)u+
 \frac{\psi_2}{\varepsilon_2}\right].
\]
Переключающая функция
\[
 \pi(\psi)=\frac{\psi_1}{\varepsilon_1}-\frac{\psi_2}{\varepsilon_2}.
\]
Отсюда
\[
 u^*=\begin{cases}
 1,& \pi>0,\\
 0,& \pi<0,\\
 [0,1],& \pi=0.
 \end{cases}
\]

Особый режим возникает, когда \(\pi(t)\equiv0\) на интервале. Из условий \(\pi=0\) и \(\dot\pi=0\) получается
\[
 y_1(t)=y_2(t),\qquad u_{\mathrm{s}}=\varepsilon_1.
\]
Следовательно, особый луч:
\[
 L_{\mathrm{s}}=\{y\in\R_+^2:\ y_1=y_2\}.
\]

Качественный вид решения при достаточно большом \(T\):
\begin{itemize}
 \item если \(y^0\in L_{\mathrm{s}}\), то на особом участке \(u=\varepsilon_1\), а в конце \(u=0\);
 \item если \(y^0\) ниже особого луча, то сначала \(u=0\), затем \(u=\varepsilon_1\), затем \(u=0\);
 \item если \(y^0\) выше особого луча, то сначала \(u=1\), затем \(u=\varepsilon_1\), затем \(u=0\).
\end{itemize}
Экзаменационная формулировка: оптимальная траектория сначала выводит систему на особый луч \(y_1=y_2\), затем некоторое время движется по нему с управлением \(u=\varepsilon_1\), а в конце весь выпуск направляется во второй сектор, поскольку функционал терминальный и равен \(y_2(T)\).

\ticket{6}{Принцип максимума Понтрягина на бесконечном интервале}

Рассматривается задача
\[
 \dot x(t)=f(x(t),u(t)),\quad x(0)=x_0,\quad u(t)\in U,
\]
\[
 J(x,u)=\int_0^\infty e^{-\rho t}g(x(t),u(t))\dd t\to\max.
\]
Для существования оптимального процесса обычно предполагают ограниченность роста системы, выпуклость множества скоростей и доходов, а также сходимость хвоста интеграла.

Функция Гамильтона--Понтрягина:
\[
 H(x,t,u,\psi,\psi^0)=\langle f(x,u),\psi\rangle+\psi^0e^{-\rho t}g(x,u),
\]
где \(\psi(t)\in\R^n\), \(\psi^0\ge0\), и пара \((\psi,\psi^0)\) нетривиальна:
\[
 \|\psi(0)\|+\psi^0>0.
\]

Если \((x^*,u^*)\) оптимальна, то существует нетривиальная сопряженная пара, для которой выполняются основные соотношения ПМП:
\[
 \dot x^*(t)=f(x^*(t),u^*(t)),
\]
\[
 \dot\psi(t)=-H_x(x^*(t),t,u^*(t),\psi(t),\psi^0),
\]
\[
 H(x^*(t),t,u^*(t),\psi(t),\psi^0)
 =\max_{u\in U}H(x^*(t),t,u,\psi(t),\psi^0)
\]
почти всюду по \(t\ge0\).

В нормальном случае можно положить \(\psi^0=1\). Тогда условие максимума сравнивает мгновенный дисконтированный выигрыш и влияние управления на будущую траекторию.

Для бесконечного горизонта важны условия трансверсальности на бесконечности. Наиболее употребительные формы:
\[
 \lim_{t\to\infty}\psi(t)=0,
\qquad\text{или}\qquad
 \lim_{t\to\infty}\langle x^*(t),\psi(t)\rangle=0.
\]
Они исключают ситуацию, когда в бесконечно далеком будущем остается неучтенная стоимость состояния.

Итог: ПМП на бесконечном интервале имеет ту же структуру, что и на конечном, но требует сходимости функционала и специальных условий трансверсальности при \(t\to\infty\).

\ticket{7}{Метод Беллмана и ПМП на бесконечном интервале. Текущие сопряженные переменные}

Пусть задана задача
\[
 \dot x=f(x,u),\quad x(0)=x_0,\quad
 J=\int_0^\infty e^{-\rho t}g(x(t),u(t))\dd t\to\max.
\]
Функция цены \(V(x)\) равна максимальному будущему дисконтированному доходу, если система начинается из состояния \(x\).
Для стационарной задачи уравнение Беллмана имеет вид
\[
 \rho V(x)=\max_{u\in U}\{g(x,u)+\langle V_x(x),f(x,u)\rangle\}.
\]
Если существует гладкое решение \(V\) и управление
\[
 u^*(x)\in\argmax_{u\in U}\{g(x,u)+\langle V_x(x),f(x,u)\rangle\},
\]
то обратная связь \(u(t)=u^*(x(t))\) является оптимальной.

Связь с принципом максимума дается через текущую сопряженную переменную:
\[
 p(t)=V_x(x^*(t)),\qquad \psi(t)=e^{-\rho t}p(t).
\]
Обычный гамильтониан в текущих ценах:
\[
 M(x,u,p)=g(x,u)+\langle p,f(x,u)\rangle.
\]
Тогда условия ПМП принимают форму
\[
 u^*(t)\in\argmax_{u\in U}M(x^*(t),u,p(t)),
\]
\[
 \dot p(t)=\rho p(t)-M_x(x^*(t),u^*(t),p(t)).
\]
Уравнение Беллмана на оптимальной траектории дает
\[
 \rho V(x^*(t))=M(x^*(t),u^*(t),p(t)).
\]

Экономический смысл:
\begin{itemize}
 \item \(V(x)\) -- внутренняя стоимость запаса капитала или состояния системы;
 \item \(V_x(x)\) -- предельная стоимость единицы состояния;
 \item \(p(t)\) -- текущие сопряженные цены;
 \item \(\psi(t)=e^{-\rho t}p(t)\) -- приведенные к начальному моменту сопряженные цены.
\end{itemize}
Условие максимума означает, что оптимальное управление максимизирует сумму текущей прибыли и прироста внутренней стоимости состояния.
Условие трансверсальности
\[
 \lim_{t\to\infty}e^{-\rho t}p(t)=0
\]
означает, что приведенная предельная стоимость капитала на бесконечности исчезает.

\ticket{8}{Достаточные условия оптимальности на бесконечном интервале}

Для задачи
\[
 \dot x=f(x,u),\quad x(0)=x_0,\quad
 J=\int_0^\infty e^{-\rho t}g(x,u)\dd t\to\max
\]
необходимые условия ПМП сами по себе не гарантируют оптимальность. Достаточные условия обычно получают из вогнутости или из метода Беллмана.

Первый вариант -- достаточность по Понтрягину. Пусть существует допустимый процесс \((x^*,u^*)\), сопряженная переменная \(\psi\), \(\psi^0=1\), и выполнены:
\[
 \dot\psi=-H_x(x^*,t,u^*,\psi,1),
\]
\[
 H(x^*,t,u^*,\psi,1)=\max_{u\in U}H(x^*,t,u,\psi,1).
\]
Если для каждого \(t\) функция
\[
 x\mapsto \max_{u\in U}H(x,t,u,\psi(t),1)
\]
вогнута по \(x\), а также выполнено условие трансверсальности
\[
 \limsup_{T\to\infty}\langle \psi(T),x(T)-x^*(T)\rangle\ge0
\]
для любой допустимой траектории \(x\), то процесс \((x^*,u^*)\) оптимален.

Второй вариант -- проверочная теорема Беллмана. Если существует гладкая функция \(V\), удовлетворяющая
\[
 \rho V(x)=\max_{u\in U}\{g(x,u)+\langle V_x(x),f(x,u)\rangle\},
\]
и управление \(u^*(x)\) достигает максимума, то соответствующий процесс оптимален, если интегральный функционал сходится и выполнено предельное условие
\[
 \lim_{T\to\infty}e^{-\rho T}V(x(T))=0
\]
для рассматриваемых допустимых процессов.

Идея доказательства в обоих случаях одна: сравнить произвольный допустимый процесс с кандидатом, проинтегрировать производную выражения \(\langle\psi,x-x^*\rangle\) или \(e^{-\rho t}V(x(t))\), а затем воспользоваться вогнутостью и предельным условием.

Итоговая формулировка: на бесконечном интервале достаточность требует не только максимума гамильтониана, но и глобального свойства вогнутости либо функции Беллмана, а также условия на поведение состояния и сопряженных переменных при \(t\to\infty\).

\ticket{9}{Геометрическая теория управления: многообразия, поля, скобки Ли}

Геометрическая теория управления изучает управляемые системы на гладких многообразиях. Гладкое многообразие \(M\) -- пространство, локально устроенное как \(\R^n\), с гладкими переходами между координатными картами.

Векторное поле \(X\) на \(M\) сопоставляет каждой точке \(q\in M\) касательный вектор \(X(q)\in T_qM\).
Обыкновенное дифференциальное уравнение
\[
 \dot q=X(q)
\]
задает поток, то есть движение точки по интегральным кривым поля.

Коммутатор двух векторных полей \(X,Y\) определяется как
\[
 [X,Y](\varphi)=X(Y\varphi)-Y(X\varphi)
\]
для гладких функций \(\varphi\). В координатах, если
\[
 X=\sum_i X_i\frac{\partial}{\partial x_i},\qquad
 Y=\sum_i Y_i\frac{\partial}{\partial x_i},
\]
то
\[
 [X,Y]=\sum_i\left(\sum_j X_j\frac{\partial Y_i}{\partial x_j}
 -Y_j\frac{\partial X_i}{\partial x_j}\right)\frac{\partial}{\partial x_i}.
\]
Линейная оболочка полей и всех их повторных коммутаторов образует алгебру Ли.

Для системы, линейной по управлению,
\[
 \dot q=f_0(q)+\sum_{i=1}^m u_i f_i(q),
\]
поля \(f_i\) задают возможные мгновенные направления движения, а их коммутаторы описывают направления, достижимые за счет быстрых чередований управлений.

Условие скобочной порождаемости:
\[
 \operatorname{Lie}\{f_1,\dots,f_m\}(q)=T_qM
\]
в каждой точке \(q\). Если оно выполнено для системы без дрейфа
\[
 \dot q=\sum_{i=1}^m u_i f_i(q),
\]
то по теореме Чоу--Рашевского любые две близкие точки одной связной компоненты можно соединить допустимой траекторией.

Экзаменационный смысл: геометрическая теория переводит вопрос управляемости в вопрос о ранге системы векторных полей и их скобок Ли.
Если исходных направлений мало, но их коммутаторы порождают все касательное пространство, система все равно может быть управляема.

\ticket{10}{Сопряженная точка и достаточные условия слабого минимума в вариационном исчислении}

Рассматривается классическая задача вариационного исчисления:
\[
 J[x]=\int_{t_0}^{t_1}L(t,x,\dot x)\dd t\to\min,\qquad
 x(t_0)=x_0,\quad x(t_1)=x_1.
\]
Необходимое условие экстремума -- уравнение Эйлера--Лагранжа:
\[
 \frac{\dd}{\dd t}L_{\dot x}(t,x,\dot x)-L_x(t,x,\dot x)=0.
\]
Экстремаль -- траектория, удовлетворяющая этому уравнению и граничным условиям.

Для исследования минимума рассматривают вторую вариацию:
\[
 \delta^2J[h]=\int_{t_0}^{t_1}
 \left(\langle P(t)\dot h,\dot h\rangle
 +2\langle Q(t)h,\dot h\rangle
 +\langle R(t)h,h\rangle\right)\dd t,
\]
где \(h(t_0)=h(t_1)=0\).
Условие Лежандра:
\[
 P(t)=L_{\dot x\dot x}(t,x^*,\dot x^*)\ge0,
\]
а усиленное условие Лежандра требует \(P(t)>0\).

Сопряженная точка к \(t_0\) -- такая точка \(\tau\in(t_0,t_1]\), для которой существует нетривиальное решение уравнения Якоби, обращающееся в нуль при \(t_0\) и \(\tau\).
Уравнение Якоби является уравнением для вариаций экстремали и получается из второй вариации.

Достаточное условие слабого минимума: если экстремаль \(x^*\) удовлетворяет уравнению Эйлера--Лагранжа, усиленному условию Лежандра и на интервале \((t_0,t_1]\) нет сопряженных к \(t_0\) точек, то \(x^*\) дает строгий слабый минимум.

Смысл условия: отсутствие сопряженных точек означает положительную определенность второй вариации на допустимых вариациях.
Поэтому малые отклонения от экстремали увеличивают значение функционала.

\ticket{11}{Особые траектории. Метод построения особых управлений}

Рассмотрим управляемую систему
\[
 \dot x=f(x,u),\qquad u\in U,
\]
и гамильтониан Понтрягина
\[
 H(x,\psi,u)=\langle \psi,f(x,u)\rangle.
\]
Если управление входит линейно,
\[
 \dot x=f_0(x)+\sum_{i=1}^m u_i f_i(x),
\]
то
\[
 H=h_0(x,\psi)+\sum_{i=1}^m u_i h_i(x,\psi),
\qquad h_i=\langle\psi,f_i(x)\rangle.
\]
Функции \(h_i\) называются переключающими.

Обычный случай: знак переключающей функции определяет крайнее управление.
Особый режим возникает, если на некотором интервале
\[
 h_i(t)\equiv0
\]
для одной или нескольких компонент управления. Тогда условие максимума не определяет управление непосредственно.

Метод построения особого управления:
\begin{enumerate}
 \item записать гамильтониан и переключающую функцию \(\Phi(t)=H_u(x,\psi,u)\);
 \item положить \(\Phi(t)\equiv0\) на особом интервале;
 \item последовательно дифференцировать:
 \[
 \Phi=0,\quad \dot\Phi=0,\quad \ddot\Phi=0,\dots;
 \]
 \item продолжать до тех пор, пока в производной явно не появится управление \(u\);
 \item решить полученное уравнение относительно \(u_{\mathrm{s}}\);
 \item проверить допустимость \(u_{\mathrm{s}}\in U\) и условие максимума.
\end{enumerate}

В терминах скобок Пуассона:
\[
 \dot h_i=\{H,h_i\}.
\]
Поэтому особые условия можно записывать через последовательные скобки гамильтонианов.

Если \(\Phi=H_u\), то порядок особого управления \(q\) определяется тем, что
\[
 \frac{\partial}{\partial u}
 \left(\frac{\dd}{\dd t}\right)^k
 \frac{\partial H}{\partial u}=0,\qquad k=0,\dots,2q-1,
\]
но
\[
 \frac{\partial}{\partial u}
 \left(\frac{\dd}{\dd t}\right)^{2q}
 \frac{\partial H}{\partial u}\ne0.
\]
Условие Келли для максимума:
\[
 (-1)^q
 \frac{\partial}{\partial u}
 \left(\frac{\dd}{\dd t}\right)^{2q}
 \frac{\partial H}{\partial u}<0.
\]
Если условие Келли нарушается или порядок приводит к частым переключениям, на соседних неособых участках может возникать chattering.

Особая траектория -- траектория, соответствующая особому управлению.
Она часто соединяется с участками bang-bang управления, то есть с участками, где управление принимает граничные значения.
Экзаменационный итог: особый режим не угадывается, а выводится из тождества переключающей функции, ее производных и условия Келли.

\ticket{12}{Граничное управление и наблюдение для волнового уравнения. Двойственность}

Типичная задача граничного управления для волнового уравнения:
\[
\begin{cases}
 y_{tt}=y_{xx}-q(x)y,\quad 0<t<T,\ 0<x<l,\\
 y(t,0)=u_0(t),\quad y(t,l)=u_1(t),\\
 y(0,x)=y_t(0,x)=0.
\end{cases}
\]
Требуется подобрать управление \(u=(u_0,u_1)\), чтобы в момент \(T\)
\[
 y(T,x)=f^0(x),\qquad y_t(T,x)=f^1(x).
\]
Это задача точной управляемости.

Естественные энергетические пространства:
\[
 y(T)\in H_0^1(0,l),\qquad y_t(T)\in L_2(0,l),
\]
а сопряженное пространство состояний можно записывать как
\[
F=H_0^1(0,l)\times L_2(0,l),\qquad
F^*=L_2(0,l)\times H^{-1}(0,l).
\]
В операторной форме
\[
 A:u\mapsto (y(T),y_t(T)),\qquad Au=f.
\]
Точная управляемость означает \(\operatorname{Im}A=F\), то есть достижимость любого конечного состояния \(f\in F\).

Сопряженная задача наблюдения:
\[
\begin{cases}
 p_{tt}=p_{xx}-q(x)p,\\
 p(t,0)=p(t,l)=0,\\
 p(T,x)=v^0(x),\quad p_t(T,x)=-v^1(x).
\end{cases}
\]
Наблюдаемым сигналом являются граничные производные:
\[
 g_0(t)=p_x(t,0),\qquad g_1(t)=-p_x(t,l).
\]

Двойственность состоит в том, что оператор управления \(A:u\mapsto (y(T),y_t(T))\) и оператор наблюдения \(A^*\) связаны интегральным тождеством:
\[
 \langle Au,v\rangle_F=\langle u,A^*v\rangle_H.
\]
Поэтому свойства управляемости исходной системы эквивалентны свойствам наблюдаемости сопряженной.

Критическое время \(T_*\) -- минимальное время, за которое волна успевает передать информацию от границы ко всей струне.
Для простейшей струны длины \(l\) при двустороннем наблюдении типично \(T_*=l\), а при одностороннем -- \(T_*=2l\).
Если \(T>T_*\), выполняется неравенство наблюдаемости, и система точно управляема; если \(T<T_*\), точная управляемость нарушается.
Экзаменационный итог: управление задается оператором \(A\), наблюдение -- сопряженным оператором \(A^*\), а ключ к точной управляемости дает неравенство наблюдаемости после критического времени.

\ticket{13}{Неравенство наблюдаемости, истокопредставимость и вариационный метод}

Пусть \(A:H\to F\) -- линейный ограниченный оператор между гильбертовыми пространствами.
Задача управления записывается как
\[
 Au=f.
\]
Система вполне управляема, если для любого \(f\in F\) уравнение имеет решение.
Сопряженная задача наблюдения:
\[
 A^*v=g.
\]
Система наблюдаема, если состояние \(v\) однозначно определяется по наблюдаемому сигналу \(A^*v\).

Для линейного ограниченного оператора в гильбертовых пространствах используется разложение
\[
 F=\overline{\operatorname{Im}A}\oplus\operatorname{Ker}A^*.
\]
Поэтому полная управляемость связана с отсутствием ненаблюдаемых состояний сопряженной системы.

Критерий управляемости через наблюдаемость:
\[
 \operatorname{Im}A=F
 \quad\Longleftrightarrow\quad
 \exists \mu>0:\ \|A^*v\|_{H^*}^2\ge \mu\|v\|_{F^*}^2
 \quad \forall v\in F^*.
\]
Это и есть неравенство наблюдаемости.
Оно означает, что наблюдаемый сигнал не может быть малым для ненулевого состояния сопряженной системы.

Нормальное решение задачи управления -- решение минимальной нормы:
\[
 u_*=\argmin\{\|u\|_H:\ Au=f\}.
\]
При неравенстве наблюдаемости оно истокопредставимо:
\[
 u_*=J_HA^*v_*,
\]
где \(J_H:H^*\to H\) -- оператор Рисса, а \(v_*\in F^*\).

Вариационный метод ищет \(v_*\) как минимум функционала
\[
 I(v)=\|A^*v\|_{H^*}^2-2\langle v,f\rangle_{F^*,F}
\]
на шаре \(\|v\|_{F^*}\le r\), если известно, что \(\|v_*\|\le r\).
После нахождения \(v_*\) управление восстанавливается по формуле
\[
 u_*=J_HA^*v_*.
\]

Приближенный метод заменяет \(A,A^*,f\) конечномерными согласованными приближениями \(\widetilde A,\widetilde A^*,\tilde f\) и решает задачу квадратичной минимизации
\[
 \widetilde I(v)=\|\widetilde A^*v\|^2-2\langle v,\tilde f\rangle\to\min.
\]
Если приближения согласованы, а ошибка данных и точность минимизации стремятся к нулю, то
\[
 \|\tilde u-u_*\|_H\to0.
\]
Условия применимости: нормальное решение должно быть истокопредставимо, известен радиус \(\|v_*\|\le r\), операторы \(A,A^*\) и данные \(f\) аппроксимируются согласованно, а ошибки данных и минимизации стремятся к нулю.
Экзаменационный итог: неравенство наблюдаемости превращает управляемость в устойчивую задачу минимизации для сопряженной переменной \(v\).

\ticket{14}{Неравенство наблюдаемости для струны и конечномерная реализация вариационного метода}

Рассматривается волновое уравнение для однородной струны:
\[
\begin{cases}
 y_{tt}=y_{xx}-q(x)y,\quad 0<t<T,\ 0<x<l,\\
 y(t,0)=u_0(t),\quad y(t,l)=u_1(t),\\
 y(0,x)=y_t(0,x)=0.
\end{cases}
\]
Цель управления:
\[
 y(T,x)=f^0(x),\qquad y_t(T,x)=f^1(x).
\]
Сопряженная задача:
\[
\begin{cases}
 p_{tt}=p_{xx}-q(x)p,\\
 p(t,0)=p(t,l)=0,\\
 p(T,x)=v^0(x),\quad p_t(T,x)=-v^1(x).
\end{cases}
\]
Наблюдение задается через
\[
 A^*v=(p_x(t,0),-p_x(t,l)).
\]

Для простейшего случая \(q(x)=0\) энергетическое тождество для сопряженной системы дает оценку
\[
 \|A^*v\|^2\ge \mu\|v\|^2,\qquad
 \mu=\frac{2(T-l)}{l}>0,\quad T>l.
\]
Следовательно, при \(T>l\) выполнено неравенство наблюдаемости, а исходная граничная система точно управляема.

Для численной реализации вводят сетку \(x_i=ih\), \(t_j=j\tau\) и аппроксимируют волновое уравнение разностной схемой:
\[
 \frac{y_i^{j+1}-2y_i^j+y_i^{j-1}}{\tau^2}
 =
 \frac{y_{i+1}^j-2y_i^j+y_{i-1}^j}{h^2}-q_i y_i^j.
\]
Граничные значения \(y_0^j=u_0^j\), \(y_N^j=u_1^j\) играют роль дискретных управлений.

Строятся взаимно сопряженные конечномерные операторы \(\widetilde A\) и \(\widetilde A^*\), чтобы сохранялась дискретная двойственность:
\[
 \langle \widetilde Au,v\rangle=\langle u,\widetilde A^*v\rangle.
\]
После этого нормальное управление ищется в виде
\[
 \tilde u=J_H\widetilde A^*v,
\]
где \(v\) минимизирует конечномерный квадратичный функционал
\[
 \|\widetilde A^*v\|^2-2\langle f,v\rangle\to\min,\qquad \|v\|\le r.
\]
Главное требование к конечномерной реализации -- сохранять двойственность операторов управления и наблюдения.
Экзаменационный итог: дискретизация должна сохранять сопряженность операторов; тогда вариационный метод дает устойчивое приближение нормального управления.

\ticket{15}{Обобщенная производная по Соболеву. Пространства Соболева. Теоремы вложения}

Пусть \(f\in L_1(a,b)\). Функция \(\omega\in L_1(a,b)\) называется обобщенной производной функции \(f\), если
\[
 \int_a^b f(x)\varphi'(x)\dd x
 =
 -\int_a^b \omega(x)\varphi(x)\dd x
 \quad \forall \varphi\in C_0^\infty(a,b).
\]
Пишут \(\omega=f'\) в обобщенном смысле.
Это определение переносит формулу интегрирования по частям на негладкие функции.

Основные свойства:
\begin{itemize}
 \item обобщенная производная единственна с точностью до равенства почти всюду;
 \item операция линейна;
 \item если классическая производная существует и интегрируема, то она совпадает с обобщенной;
 \item если обобщенная производная равна нулю, то функция постоянна почти всюду на связном интервале.
\end{itemize}

Пространство Соболева
\[
 W_p^k(\Omega)=\{u\in L_p(\Omega): D^\alpha u\in L_p(\Omega),\ |\alpha|\le k\}
\]
с нормой
\[
 \|u\|_{W_p^k}=
 \left(\sum_{|\alpha|\le k}\|D^\alpha u\|_{L_p}^p\right)^{1/p}
\]
при \(1\le p<\infty\).
При \(p=2\) пишут
\[
 H^k(\Omega)=W_2^k(\Omega).
\]
Например,
\[
 H^1(a,b)=\{u\in L_2(a,b):u'\in L_2(a,b)\},
\]
\[
 \|u\|_{H^1}^2=\int_a^b (u^2+(u')^2)\dd x.
\]

Теоремы вложения показывают, что наличие обобщенных производных повышает регулярность функции.
В одномерном случае
\[
 H^1(a,b)\subset C[a,b],
\]
то есть каждый класс из \(H^1(a,b)\) имеет непрерывного представителя, и
\[
 \|u\|_{C[a,b]}\le C\|u\|_{H^1(a,b)}.
\]
Кроме того,
\[
 u(\beta)-u(\alpha)=\int_\alpha^\beta u'(x)\dd x.
\]
Также \(C^\infty[a,b]\) плотно в \(H^1(a,b)\), то есть любую функцию Соболева можно приблизить гладкими функциями в норме \(H^1\).

Экзаменационный смысл: пространства Соболева позволяют корректно формулировать краевые задачи для уравнений в частных производных, даже когда классических производных у решения нет.

\ticket{16}{Оптимальный нагрев стержня. Обобщенное решение и теорема Вейерштрасса}

Задача оптимального нагрева стержня:
\[
 J(u)=\int_0^l (x(s,T,u)-b(s))^2\dd s\to\inf,
\]
\[
\begin{cases}
 x_t=x_{ss}+u(s,t),\qquad (s,t)\in Q=(0,l)\times(0,T),\\
 x(0,t)=x(l,t)=0,\\
 x(s,0)=\varphi(s),\\
 u\in U=\{u\in L_2(Q):\|u\|_{L_2(Q)}\le R\}.
\end{cases}
\]
Фазовая переменная \(x\) -- температура, управление \(u\) -- распределенный нагрев, цель -- приблизить конечное распределение температуры к заданной функции \(b(s)\).

Обобщенное решение \(x(s,t)\) определяется через интегральное тождество. Требуется
\[
 x\in L_2(Q),\qquad x_s\in L_2(Q),\qquad x|_{s=0}=x|_{s=l}=0,
\]
и для всех допустимых пробных функций \(\Phi\)
\[
 \int_0^l x(s,T)\Phi(s,T)\dd s
 +\iint_Q\bigl(-x\Phi_t+x_s\Phi_s\bigr)\dd s\dd t
 =
 \int_0^l \varphi(s)\Phi(s,0)\dd s
 +\iint_Q u\Phi\dd s\dd t.
\]
Это слабая форма уравнения теплопроводности.

Теорема существования и единственности: для каждого \(u\in L_2(Q)\) краевая задача имеет единственное обобщенное решение.
Доказательство обычно основано на энергетических оценках и плотности гладких функций.
Единственность получается из рассмотрения разности двух решений: если \(y=x_1-x_2\), то \(y\) решает однородную задачу, энергетическая оценка дает \(y=0\).

Существование оптимального управления следует из теоремы Вейерштрасса в слабой форме.
Множество \(U\) выпукло, замкнуто и ограничено в гильбертовом пространстве \(L_2(Q)\), значит оно слабо компактно.
Если \(u_k\) -- минимизирующая последовательность, то из нее можно выделить слабо сходящуюся подпоследовательность \(u_k\rightharpoonup u_*\in U\).
Оператор \(u\mapsto x(\cdot,T,u)\) непрерывен в нужной слабой топологии, а функционал \(J\) слабо полунепрерывен снизу.
Следовательно,
\[
 J(u_*)=\inf_{u\in U}J(u),
\]
то есть оптимальный нагрев существует.

\ticket{17}{Градиент функционала в задаче нагрева. Критерий оптимальности и РМПГ}

Для задачи нагрева
\[
 J(u)=\int_0^l (x(s,T,u)-b(s))^2\dd s\to\inf,\qquad u\in U
\]
градиент функционала выражается через сопряженную задачу.
Если \(x\) -- решение прямой задачи, то сопряженная функция \(\psi\) удовлетворяет
\[
\begin{cases}
 -\psi_t=\psi_{ss},\\
 \psi(0,t)=\psi(l,t)=0,\\
 \psi(s,T)=2(x(s,T,u)-b(s)).
\end{cases}
\]
Тогда производная Фреше имеет вид
\[
 J'(u)=\psi,
\]
то есть для вариации \(v\)
\[
 \langle J'(u),v\rangle=\iint_Q \psi(s,t)v(s,t)\dd s\dd t.
\]

Если \(U=\{u\in L_2(Q):\|u\|\le R\}\), то критерий оптимальности:
\[
 u_*\in U_*
 \quad\Longleftrightarrow\quad
 \langle J'(u_*),u-u_*\rangle\ge0\quad \forall u\in U.
\]
Для шара это эквивалентно условиям Куна--Таккера:
\[
 J'(u_*)+\lambda_*u_*=0,\qquad
 \lambda_*(\|u_*\|-R)=0,\qquad
 \|u_*\|\le R,\quad \lambda_*\ge0.
\]
Если ограничение активно, то
\[
 \lambda_*=\frac{\|J'(u_*)\|}{R}.
\]

Регуляризованный метод проекции градиента строит последовательность
\[
 u_{k+1}=\pi_U\bigl(u_k-\beta_k(J_k'(u_k)+2\alpha_k u_k)\bigr),
\]
где \(\pi_U\) -- проекция на \(U\), \(\alpha_k>0\) -- параметр регуляризации Тихонова, а \(J_k'\) -- приближенный градиент.
Регуляризованный функционал:
\[
 T_k(u)=J(u)+\alpha_k\|u\|^2,\qquad T_k'(u)=J'(u)+2\alpha_k u.
\]

При выпуклости \(J\), липшицевом типе условия на градиент, согласовании \(\alpha_k,\beta_k\) и ошибки приближенного градиента \(\delta_k\) метод сходится:
\[
 J(u_k)\to J_*,\qquad u_k\to u_*
\]
к нормальному оптимальному управлению.

\ticket{18}{Равновесие по Нэшу в линейно-квадратичной дифференциальной игре трех лиц}

Рассматривается линейно-квадратичная дифференциальная игра с двумя игроками и неопределенностью, которая трактуется как третий игрок:
\[
 \dot x=A(t)x+u_1+u_2+z,\qquad x(t_*)=x_*.
\]
Стратегии имеют линейную позиционную форму
\[
 u_i(t,x)=Q_i(t)x,\qquad z(t,x)=P(t)x.
\]
Функционалы первых двух игроков квадратичны:
\[
 J_i=x^\top(\theta)C_i x(\theta)
 +\int_{t_*}^{\theta}
 \left(\sum_{j=1}^2 u_j^\top D_{ij}u_j+z^\top L_i z+x^\top G_i x\right)\dd t,
 \quad i=1,2.
\]
Третий игрок задается как
\[
 J_3=-\alpha J_1-(1-\alpha)J_2,\qquad \alpha\in(0,1).
\]

Тройка \((U_1^e,U_2^e,Z^G)\) является равновесием по Нэшу, если ни один участник не может увеличить свой функционал односторонним отклонением:
\[
 J_1(U_1,U_2^e,Z^G)\le J_1(U_1^e,U_2^e,Z^G),
\]
\[
 J_2(U_1^e,U_2,Z^G)\le J_2(U_1^e,U_2^e,Z^G),
\]
\[
 J_3(U_1^e,U_2^e,Z)\le J_3(U_1^e,U_2^e,Z^G).
\]

Достаточное условие существования формулируется через функции Беллмана \(V_i(t,x)\).
Если существуют функции \(V_1,V_2,V_3\) и максимизаторы
\[
 u_1(t,x,V),\quad u_2(t,x,V),\quad z(t,x,V),
\]
которые удовлетворяют системе уравнений Гамильтона--Якоби--Беллмана
\[
 W_i(t,x,u_1,u_2,z,V_i)=0,\qquad i=1,2,3,
\]
с терминальными условиями
\[
 V_i(\theta,x)=x^\top C_i x,\quad i=1,2,\qquad
 V_3(\theta,x)=-x^\top C(\alpha)x,
\]
то стратегии
\[
 U_i^e(t,x)=u_i(t,x,V^*(t,x)),\qquad
 Z^G(t,x)=z(t,x,V^*(t,x))
\]
образуют равновесие по Нэшу.

В линейно-квадратичном случае функции Беллмана ищутся в квадратичном виде
\[
 V_i(t,x)=x^\top K_i(t)x.
\]
Тогда задача сводится к системе матричных уравнений типа Риккати, а равновесные стратегии имеют линейную форму по состоянию.
Если эта связанная система Риккати имеет решение на всем интервале, то полученные линейные обратные связи образуют равновесие по Нэшу.
Экзаменационный итог: в ЛК-игре равновесие ищется через квадратичные функции Беллмана, а проверка сводится к разрешимости системы Риккати.

\ticket{19}{Равновесие по Бержу в ЛК-игре двух лиц с малым влиянием второго игрока}

Игра задается системой
\[
 \dot x=A(t)x+u_1+\varepsilon u_2,\qquad x(t_0)=x_0,
\]
где \(0\le\varepsilon\ll1\). Игроки используют линейные стратегии
\[
 u_i(t,x)=Q_i(t)x.
\]
Функционалы имеют линейно-квадратичный вид:
\[
 J_i=x^\top(\theta)C_i x(\theta)+
 \int_{t_0}^{\theta}\left(u_1^\top D_{i1}u_1+u_2^\top D_{i2}u_2\right)\dd t,
 \quad i=1,2.
\]

Равновесие по Бержу отличается от равновесия Нэша.
В равновесии Нэша каждый игрок не может улучшить свой выигрыш своим отклонением.
В равновесии Бержу каждый игрок получает максимум благодаря выбору стратегий остальными:
\[
 \max_{U_2}J_1(U_1^B,U_2)=J_1(U^B),
\qquad
 \max_{U_1}J_2(U_1,U_2^B)=J_2(U^B).
\]

Для построения равновесия вводятся функции Беллмана \(V_i(t,x)\) и гамильтонианы
\[
 W_i=V_{it}+V_{ix}(Ax+u_1+\varepsilon u_2)
 +u_1^\top D_{i1}u_1+u_2^\top D_{i2}u_2.
\]
Условия максимума по стратегии другого игрока дают
\[
 u_1^B(t,x)=-D_{21}^{-1}\Theta_2(t,\varepsilon)x,
\]
\[
 u_2^B(t,x)=-\varepsilon D_{12}^{-1}\Theta_1(t,\varepsilon)x,
\]
если
\[
 V_i(t,x)=x^\top\Theta_i(t,\varepsilon)x.
\]
Матрицы \(\Theta_i\) удовлетворяют системе матричных дифференциальных уравнений типа Риккати с терминальными условиями
\[
 \Theta_i(\theta,\varepsilon)=C_i.
\]

При \(\varepsilon=0\) система упрощается; если она имеет продолжимое решение на \([t_0,\theta]\), то по теореме Пуанкаре при достаточно малых \(\varepsilon\) существует продолжимое решение и для исходной игры.

Экзаменационный итог: \(\varepsilon=0\) является базовой задачей, а при малом \(\varepsilon\) равновесие по Бержу строится как возмущение этого решения через уравнения Риккати.

\ticket{20}{Максимум по Парето в двухкритериальной линейно-квадратичной динамической задаче}

Пусть лицо, принимающее решение, выбирает стратегию \(x\in X\subset\R^n\), а неопределенность \(y\in Y\subset\R^m\) заранее неизвестна.
Есть два критерия: выигрыш \(f(x,y)\) и отрицательный риск \(-R(x,y)\).
Задача записывается как
\[
 G_2=\langle X,Y,\{f(x,y),-R(x,y)\}\rangle.
\]
Функция выигрыша линейно-квадратична:
\[
 f(x,y)=x^\top Ax+2x^\top By+y^\top Dy+2a^\top x+2d^\top y,
\]
где \(A,D\) симметричны.

Функция риска или сожаления:
\[
 R(x,y)=\max_{z\in X}f(z,y)-f(x,y).
\]
Она показывает, сколько выигрыша теряется из-за выбора \(x\), если реализовалась неопределенность \(y\).

Неопределенность \(y^P(x)\) называется минимальной по Парето для фиксированного \(x\), если нельзя одновременно увеличить выигрыш и уменьшить риск:
\[
 f(x,y)\le f(x,y^P),\qquad -R(x,y)\le -R(x,y^P),
\]
причем хотя бы одно неравенство строгое, невозможны одновременно.

Достаточная конструкция: выбирается \(\sigma\in(0,1)\) и функция
\[
 \Phi(x,y)=f(x,y)-\sigma f[y],
\qquad
 f[y]=\max_{x\in X}f(x,y).
\]
Если
\[
 y^P(x)\in\argmin_{y\in Y}\Phi(x,y),
\]
то \(y^P(x)\) является минимальной по Парето неопределенностью.

Стратегия \(x^P\) максимальна по Парето, если в задаче
\[
 \langle X,\{f(x,y^P(x)),-R(x,y^P(x))\}\rangle
\]
не существует другой стратегии \(x\), которая не хуже по обоим критериям и лучше хотя бы по одному.
Практически \(x^P\) часто ищут из условия максимизации скаляризованной функции \(\Phi(x,y^P(x))\).

Парето-максимум вообще говоря не единственен: это множество стратегий, которые нельзя улучшить сразу по всем критериям.
Экзаменационный итог: Парето-подход заменяет один критерий множеством неулучшаемых компромиссов между выигрышем и риском; на практике используют скаляризацию.

\ticket{21}{Необходимые и достаточные условия существования стратегий с нулевым риском}

Рассматривается задача принятия решений при неопределенности
\[
 \langle X,Y,f(x,y)\rangle,
\]
где \(x\in X\) -- стратегия, \(y\in Y\) -- неопределенность, \(f(x,y)\) -- выигрыш.
Функция риска:
\[
 \Phi(x,y)=\max_{z\in X}f(z,y)-f(x,y).
\]
Она неотрицательна:
\[
 \Phi(x,y)\ge0.
\]
Риск равен нулю тогда и только тогда, когда выбранная стратегия \(x\) дает максимальный выигрыш при данной неопределенности \(y\).

Если \(X,Y\) компакты, а \(f\) непрерывна, то \(\Phi\) непрерывна, и минимаксная задача по риску имеет решение:
\[
 \min_{x\in X}\max_{y\in Y}\Phi(x,y).
\]

Стратегия \(x^r\) имеет нулевой риск при всех неопределенностях, если
\[
 \Phi(x^r,y)=0\qquad \forall y\in Y.
\]
Необходимое и достаточное условие существования такой стратегии: функция риска имеет седловую точку \((x^r,y^0)\):
\[
 \max_{y\in Y}\Phi(x^r,y)=
 \Phi(x^r,y^0)=
 \min_{x\in X}\Phi(x,y^0).
\]
Эквивалентно:
\[
 \min_{x\in X}\max_{y\in Y}\Phi(x,y)
 =
 \max_{y\in Y}\min_{x\in X}\Phi(x,y)
 =
 0.
\]

Доказательство по смыслу простое.
Если \(x^r\) дает нулевой риск для всех \(y\), то максимум риска по \(y\) равен нулю, а так как риск неотрицателен, это минимально возможное значение; значит возникает седловая точка.
Обратно, если седловая точка имеет значение ноль, то максимальный риск стратегии \(x^r\) равен нулю, следовательно \(\Phi(x^r,y)=0\) для всех \(y\).

Экзаменационный итог: нулевой риск -- очень сильное свойство; оно означает, что одна стратегия одновременно оптимальна для всех возможных неопределенностей.

\ticket{22}{Гарантированное по Парето равновесие в многошаговой дуополии Курно с импортом}

Рассматривается многошаговая бескоалиционная игра двух фирм при неопределенности импорта.
Состояния \(q_1(k),q_2(k)\) -- объемы выпуска фирм.
Управления \(u_1,u_2\) -- дополнительные решения игроков, а \(z(k)\) -- неопределенность, связанная с импортом.
Типичная динамика:
\[
 q_1(k+1)=\frac{a-c}{2b}-\frac{q_2(k)}2-\frac{z(k)}2+u_1(k,q(k)),
\]
\[
 q_2(k+1)=\frac{a-c}{2b}-\frac{q_1(k)}2-\frac{z(k)}2+u_2(k,q(k)).
\]
Функции выигрыша:
\[
 J_i=q_i^2(2)+\sum_{k=0}^{1}
 \left(q_i(k)-2u_i^2(k)+\frac{z^2(k)}2\right),
 \qquad i=1,2.
\]

Гарантированное по Парето равновесие -- это тройка
\[
 (U^e,J_1^e(q_0),J_2^e(q_0)),
\]
для которой существует неопределенность \(Z^P\), минимальная по Парето, такая что:
\begin{enumerate}
 \item \(Z^P\) минимальна по Парето в двухкритериальной задаче по выигрышам игроков;
 \item \(U^e\) является равновесием по Нэшу в игре, где неопределенность зафиксирована как \(Z=Z^P\);
 \item \(J_i^e=J_i(U^e,Z^P,q_0)\) -- гарантированные выигрыши игроков.
\end{enumerate}

Алгоритм построения основан на динамическом программировании назад по времени.
На последнем шаге задаются терминальные функции:
\[
 V^{(2)}(q)=q_1^2+q_2^2,\qquad V_i^{(2)}(q)=q_i^2.
\]
На шаге \(k=1\) сначала находится Парето-минимальная неопределенность \(z^P(1,q,u)\), затем при этой неопределенности находятся равновесные управления \(u_i^e(1,q)\).
Далее строятся функции \(V_i^{(1)}(q)\).
На шаге \(k=0\) та же процедура повторяется с учетом уже построенных функций продолжения.

В результате получаются:
\[
 U^e=(U_1^e,U_2^e),\qquad
 J_i^e(q_0)=V_i^{(0)}(q_0).
\]
Экзаменационный итог: порядок построения фиксирован -- сначала выбирается Парето-минимальная неопределенность импорта, затем при ней строится равновесие Нэша методом обратной индукции.

\ticket{23}{Гарантированное по выигрышам и рискам равновесие в линейно-квадратичной игре двух лиц}

Рассматривается игра
\[
 \Gamma_2=\langle\{1,2\},X_1,X_2,Y,\{f_i(x,y)\}_{i=1}^2\rangle,
\]
где \(X_1=X_2=\R^n\), \(Y=\R^m\), а выигрыши линейно-квадратичны:
\[
 f_1=x_1^\top A_1x_1+2x_1^\top x_2+y^\top D_1y+2a_1^\top x_1+2d_1^\top y,
\]
\[
 f_2=x_2^\top A_2x_2-2x_2^\top x_1+y^\top D_2y+2a_2^\top x_2+2d_2^\top y.
\]
Лучший выигрыш первого игрока при фиксированных \(x_2,y\):
\[
 f_1[x_2,y]=\max_{x_1}f_1(x_1,x_2,y),
\]
аналогично определяется \(f_2[x_1,y]\).
Риски:
\[
 R_1=f_1[x_2,y]-f_1(x_1,x_2,y),\qquad
 R_2=f_2[x_1,y]-f_2(x_1,x_2,y).
\]

Гарантированное по выигрышам и рискам равновесие -- пятерка
\[
 (x^e,f_1^P,f_2^P,R_1^P,R_2^P),
\]
для которой существует неопределенность \(y^P\), такая что
\[
 f_i^P=f_i(x^e,y^P),\qquad R_i^P=R_i(x^e,y^P),
\]
и выполняются два условия:
\begin{enumerate}
 \item \(x^e\) -- равновесие Нэша в игре при фиксированном \(y=y^P\);
 \item \(y^P\) минимальна по Парето в четырехкритериальной задаче
 \[
 \{f_1,-R_1,f_2,-R_2\}.
 \]
\end{enumerate}

Если \(A_i<0\), то максимумы по собственным стратегиям находятся из первых производных:
\[
 f_1[x_2,y]=-(x_2+a_1)^\top A_1^{-1}(x_2+a_1)+y^\top D_1y+2d_1^\top y,
\]
\[
 f_2[x_1,y]=-(x_1-a_2)^\top A_2^{-1}(x_1-a_2)+y^\top D_2y+2d_2^\top y.
\]
При \(D_i>0\) Парето-гарантированная неопределенность имеет вид
\[
 y^P=-(D_1+D_2)^{-1}(d_1+d_2).
\]
Равновесная ситуация определяется системой первых порядков:
\[
 x_1^e=(A_2A_1+I)^{-1}(a_2-A_2a_1),
\]
\[
 x_2^e=(A_1A_2+I)^{-1}(a_1+A_1a_2).
\]
После этого гарантированные выигрыши и риски получаются прямой подстановкой \(x^e,y^P\) в \(f_i,R_i\).
Экзаменационный итог: сначала из условий первого порядка находится равновесная ситуация \(x^e\), затем выбранная неопределенность \(y^P\) подставляется в выигрыши и риски.

\ticket{24}{Математические модели технического прогресса}

Технический прогресс может быть экзогенным и эндогенным.
Экзогенный прогресс задается извне: производственная функция зависит от времени или внешнего технологического уровня:
\[
 Q(t)=F(K(t),L(t),t),\qquad \frac{\partial F}{\partial t}\ge0.
\]
Здесь \(K\) -- капитал, \(L\) -- труд, \(Q\) -- выпуск.

В модели Солоу--Свона:
\[
 \dot K=I-\mu K,\qquad \dot L=\eta L,\qquad Q=F(K,L),
\]
часто вводят норму накопления
\[
 s(t)=\frac{I(t)}{Q(t)}.
\]
В модели Рамсея выбор между потреблением и инвестициями становится задачей оптимального управления:
\[
 \dot k=sf(k)-(\mu+\eta)k,\qquad
 \int_0^\infty e^{-\rho t}U(c(t))\dd t\to\max.
\]

Эндогенный прогресс возникает внутри модели как результат инвестиций в знания, образование, исследования и человеческий капитал.
Производственная функция включает дополнительную переменную, например человеческий капитал \(H\) или уровень знаний \(A\):
\[
 Q=F(K,L,H)\quad \text{или}\quad Q=A K^\alpha L^{1-\alpha}.
\]
Уровень знаний может удовлетворять уравнению
\[
 \dot A=\delta(1-u)L,
\]
где часть ресурсов направляется не в текущее производство, а в накопление знаний.

Пример двухсекторной модели:
\[
\begin{cases}
 \dot x_1=uF(x),\\
 \dot x_2=(1-u)F(x),\\
 J=x_2(T)\to\max,\qquad u\in[0,1],
\end{cases}
\]
где один сектор отвечает за производство или физический капитал, а другой -- за человеческий капитал или конечный результат.

Подходы к решению:
\begin{itemize}
 \item записать фазовые переменные и управление;
 \item построить функционал полезности или терминальный критерий;
 \item применить принцип максимума Понтрягина или метод Беллмана;
 \item исследовать стационарные и особые режимы, траектории сбалансированного роста.
\end{itemize}
Экзаменационный вывод: экзогенный прогресс задается вне оптимизационной задачи, а эндогенный объясняется оптимальным распределением ресурсов между производством, накоплением капитала и созданием знаний.

\ticket{25}{Оптимальная добыча невозобновляемых ресурсов}

Пусть \(R(t)\ge0\) -- запас невозобновляемого ресурса, \(E(t)\) -- интенсивность добычи, \(p(t)>0\) -- цена, \(C(E,R)\) -- себестоимость.
Динамика запаса:
\[
 \dot R(t)=-E(t),\qquad R(0)=R_0>0,\qquad R(T)=0.
\]
Ограничения:
\[
 0\le E(t)\le E_{\max}.
\]
Функционал прибыли:
\[
 J(E,T)=\int_0^T e^{-rt}\bigl(p(t)E(t)-C(E(t),R(t))\bigr)\dd t\to\max.
\]
Горизонт \(T\) может быть заданным или свободным.

Гамильтониан:
\[
 H=e^{-rt}\bigl(p(t)E-C(E,R)\bigr)-\psi E.
\]
Сопряженное уравнение:
\[
 \dot\psi=-H_R=e^{-rt}C_R(E,R).
\]
Условие максимума:
\[
 E^*(t)\in\argmax_{0\le E\le E_{\max}}
 \left[e^{-rt}(p(t)E-C(E,R^*))-\psi(t)E\right].
\]
Если максимум внутренний, то
\[
 e^{-rt}\bigl(p(t)-C_E(E^*,R^*)\bigr)=\psi(t).
\]

В простейшем случае без издержек \(C=0\) условие максимума сравнивает \(e^{-rt}p(t)\) с \(\psi(t)\).
Если \(e^{-rt}p(t)>\psi(t)\), выгодно добывать максимально: \(E=E_{\max}\).
Если \(e^{-rt}p(t)<\psi(t)\), добыча откладывается: \(E=0\).
При равенстве возможен особый режим.

Экономический смысл сопряженной переменной \(\psi\) -- теневая цена единицы ресурса в недрах.
Условие оптимальности говорит: добывать нужно тогда, когда текущая дисконтированная маржинальная прибыль не меньше ценности сохранения ресурса для будущего.
Для свободного \(T\) добавляется условие трансверсальности по конечному времени, обычно \(H(T)=0\), если конечное состояние фиксировано соответствующим образом.

\ticket{26}{Интегрированные модели климата и экономики. Модель DICE}

Модель DICE объединяет экономический рост, выбросы, климатическую систему и функцию общественного благосостояния.
Производство задается функцией типа Кобба--Дугласа:
\[
 Q(t)=\Omega(t)A(t)K^\gamma(t)L^{1-\gamma}(t),
\]
где \(A\) -- технология, \(K\) -- капитал, \(L\) -- труд, \(\Omega(t)\) -- климатический ущерб.
Выпуск распределяется между потреблением и инвестициями:
\[
 Q=C+I,\qquad I=S(t)Q(t),\qquad C=(1-S(t))Q(t).
\]
Капитал:
\[
 \dot K=S(t)Q(t)-\delta_KK(t).
\]

Промышленные выбросы:
\[
 E(t)=(1-\mu(t))\sigma(t)A(t)K^\gamma(t)L^{1-\gamma}(t),
\]
где \(\mu(t)\in[0,1]\) -- управление сокращением выбросов, \(\sigma(t)\) -- углеродоемкость.
Концентрация парниковых газов:
\[
 \dot M=\alpha E-\delta_M(M-\widetilde M).
\]
Радиационное воздействие:
\[
 F(t)=\eta\log_2\frac{M(t)}{\widetilde M}+o(t).
\]
Температурный блок:
\[
 \dot T_U=\frac1{R_1}\left(F-\lambda T_U-\frac{R_2}{\tau_{12}}(T_U-T_L)\right),
\]
\[
 \dot T_L=\frac1{\tau_{12}}(T_U-T_L).
\]

С учетом затрат на сокращение выбросов и ущерба выпуск можно записать в форме
\[
 Q_{\mathrm{net}}=
 \frac{(1-b_1\mu^{b_2})A K^\gamma L^{1-\gamma}}
 {1+Q_1T_U^{Q_2}}.
\]
Целевой функционал -- дисконтированное общественное благосостояние:
\[
 W=\int_{T_0}^{\infty}e^{-\rho(t-T_0)}
 L(t)\ln\frac{C(t)}{L(t)}\dd t\to\max.
\]
Управления: норма сбережений \(S(t)\) и уровень сокращения выбросов \(\mu(t)\).

Подход к решению: это задача оптимального управления с несколькими фазовыми переменными \(K,M,T_U,T_L\).
Применяются принцип максимума Понтрягина, метод Беллмана или численная динамическая оптимизация.
Экзаменационный итог: модель DICE сводит климато-экономическую задачу к оптимальному управлению, где оптимум есть компромисс между текущим потреблением, инвестициями и снижением будущего климатического ущерба.

\ticket{27}{Неантагонистические дифференциальные игры. Равновесие по Нэшу на конечном интервале}

Рассматривается игра \(N\) игроков:
\[
 \dot x=f(t,x,u_1,\dots,u_N),\qquad x(t_0)=x_0,
\]
\[
 J_i=\int_{t_0}^T F_i(t,x,u_1,\dots,u_N)\dd t+q_i(x(T)),
 \qquad i=1,\dots,N.
\]
Игра неантагонистическая, потому что функционалы игроков не обязаны быть противоположными.

Программные стратегии -- это управления, зависящие только от времени и начальной позиции:
\[
 u_i=u_i(t).
\]
Набор \((u_1^*,\dots,u_N^*)\) является равновесием по Нэшу в программных стратегиях, если для каждого игрока \(i\)
\[
 J_i(u_i,u_{-i}^*)\le J_i(u_i^*,u_{-i}^*)\qquad \forall u_i.
\]
То есть ни один игрок не может улучшить свой результат, меняя только свою стратегию.

Необходимые условия имеют вид принципа максимума для каждого игрока.
Для игрока \(i\) вводится гамильтониан
\[
 H_i=F_i(t,x,u_1,\dots,u_N)+
 \langle \psi_i,f(t,x,u_1,\dots,u_N)\rangle.
\]
Тогда на равновесной траектории:
\[
 u_i^*(t)\in\argmax_{v_i\in U_i}
 H_i(t,x^*,u_1^*,\dots,u_{i-1}^*,v_i,u_{i+1}^*,\dots,u_N^*,\psi_i),
\]
\[
 \dot\psi_i=-\frac{\partial H_i}{\partial x},\qquad
 \psi_i(T)=q_{ix}(x^*(T)).
\]

Позиционные стратегии имеют вид
\[
 u_i=\varphi_i(t,x).
\]
Равновесие по Нэшу в позиционных стратегиях определяется тем же принципом одностороннего отклонения, но остальные игроки используют обратные связи \(\varphi_j^*(t,x)\).

Состоятельное позиционное равновесие усиливает понятие: равновесие должно сохраняться в любой подыгре, начинающейся из любой точки \((t,x)\).
Оно описывается системой уравнений Гамильтона--Якоби:
\[
 -V_{it}=\max_{u_i\in U_i}
 \{F_i(t,x,\varphi_1^*,\dots,u_i,\dots,\varphi_N^*)
 +V_{ix}f(t,x,\varphi_1^*,\dots,u_i,\dots,\varphi_N^*)\}.
\]
Экзаменационный итог: равновесие Нэша проверяется односторонними отклонениями, а состоятельное позиционное равновесие -- уравнениями Гамильтона--Якоби для каждой подыгры.

\ticket{28}{Дифференциальные игры с бесконечной продолжительностью}

Для игры бесконечной продолжительности:
\[
 \dot x=f(t,x,u_1,\dots,u_N),\qquad x(t_0)=x_0,
\]
\[
 J_i=\int_{t_0}^{\infty}e^{-r(t-t_0)}
 F_i(t,x,u_1,\dots,u_N)\dd t.
\]
Дисконтирование с \(r>0\) обеспечивает сходимость функционалов и уменьшает значение дальнего будущего.

В программных стратегиях равновесие по Нэшу определяется условием
\[
 J_i(u_i,u_{-i}^*)\le J_i(u_i^*,u_{-i}^*)
 \qquad \forall u_i,\quad i=1,\dots,N.
\]
Необходимые условия ПМП:
\[
 u_i^*(t)\in\argmax_{v_i\in U_i}
 \{F_i(t,x^*,u_{-i}^*,v_i)+\langle p_i,f(t,x^*,u_{-i}^*,v_i)\rangle\},
\]
\[
 \dot p_i=r p_i-\frac{\partial}{\partial x}
 \{F_i+\langle p_i,f\rangle\}.
\]
Здесь \(p_i\) -- текущая сопряженная переменная, связанная с обычной переменной формулой
\[
 \psi_i(t)=e^{-r(t-t_0)}p_i(t).
\]
На бесконечности добавляются условия трансверсальности.

В позиционных стратегиях \(u_i=\varphi_i(x)\) равновесие означает, что ни один игрок не улучшает функционал, отклоняясь от своей обратной связи при фиксированных обратных связях остальных.
Состоятельное позиционное равновесие описывается стационарной системой Беллмана:
\[
 rW_i(x)=
 \max_{u_i\in U_i}
 \{F_i(x,\varphi_1^*,\dots,u_i,\dots,\varphi_N^*)
 +W_{ix}(x)f(x,\varphi_1^*,\dots,u_i,\dots,\varphi_N^*)\}.
\]

Экзаменационный итог: главное отличие от конечного горизонта -- нет терминального условия \(V(T,x)=q(x)\), зато появляются дисконтирование, условия сходимости и трансверсальность при \(t\to\infty\).

\ticket{29}{Модель конкурентной рекламы с двумя участниками}

На рынке две фирмы. Переменная \(x(t)\in[0,1]\) -- доля рынка первой фирмы, \(1-x(t)\) -- доля второй.
Управления \(u_1,u_2\) -- расходы на рекламу.
Динамика:
\[
 \dot x=u_1\sqrt{1-x}-u_2\sqrt{x},\qquad x(0)=x_0.
\]
Функционалы:
\[
 J_1=\int_0^T\left(q_1x(t)-\frac{c_1}{2}u_1^2(t)\right)e^{-rt}\dd t
 +e^{-rT}S_1x(T),
\]
\[
 J_2=\int_0^T\left(q_2(1-x(t))-\frac{c_2}{2}u_2^2(t)\right)e^{-rt}\dd t
 +e^{-rT}S_2(1-x(T)).
\]

Для равновесия по Нэшу в программных стратегиях вводятся сопряженные переменные \(\Lambda_1,\Lambda_2\).
Условия максимума дают
\[
 u_1^*(t)=\frac{\Lambda_1(t)}{c_1}e^{rt}\sqrt{1-x^*(t)},
\]
\[
 u_2^*(t)=\frac{\Lambda_2(t)}{c_2}e^{rt}\sqrt{x^*(t)}.
\]
Сопряженные переменные удовлетворяют системе
\[
 \dot\Lambda_1=-q_1e^{-rt}
 +\Lambda_1\left(\frac{u_1^*}{2\sqrt{1-x^*}}
 +\frac{u_2^*}{2\sqrt{x^*}}\right),
\]
\[
 \dot\Lambda_2=q_2e^{-rt}
 +\Lambda_2\left(\frac{u_1^*}{2\sqrt{1-x^*}}
 +\frac{u_2^*}{2\sqrt{x^*}}\right),
\]
\[
 \Lambda_1(T)=e^{-rT}S_1,\qquad
 \Lambda_2(T)=-e^{-rT}S_2.
\]
Подстановка управлений дает замкнутую систему для \(x^*,\Lambda_1,\Lambda_2\).

Для состоятельного позиционного равновесия используются функции Беллмана:
\[
 V_1(t,x)=e^{-rt}(A_1(t)x+B_1(t)),
\]
\[
 V_2(t,x)=e^{-rt}(A_2(t)(1-x)+B_2(t)).
\]
Тогда равновесные обратные связи имеют вид
\[
 \varphi_1^*(t,x)=\frac{A_1(t)}{c_1}\sqrt{1-x},
\qquad
 \varphi_2^*(t,x)=\frac{A_2(t)}{c_2}\sqrt{x}.
\]
Коэффициенты \(A_i,B_i\) определяются из системы обыкновенных дифференциальных уравнений с терминальными условиями \(A_i(T)=S_i,\ B_i(T)=0\).

Экономический смысл: реклама первой фирмы увеличивает ее долю рынка тем сильнее, чем больше оставшийся рынок \(1-x\); реклама второй фирмы уменьшает долю первой пропорционально \(\sqrt{x}\).
Экзаменационный итог: программное равновесие получается из ПМП, а состоятельное позиционное -- из линейной по \(x\) функции Беллмана и соответствующих обратных связей.

\end{document}
